% Part: computability
% Chapter: computability-theory
% Section: total

\documentclass[../../../include/open-logic-section]{subfiles}

\begin{document}

\olfileid{cmp}{thy}{tot}
\olsection{عدم قابلية تقرير الكلية}

ومن استعمالات مبرهنة $s$-$m$-$n$ في إثبات عدم قابلية
الحساب المثال الآتي. نجعل $\fn{Tot}$ مجموعة فهارس الدوال الكلية القابلة
للحساب، فيكون
\[
\fn{Tot} = \Setabs{x}{\text{لكل $y$، تكون $\cfind{x}(y)\fdefined$}}.
\]

\begin{prop}
\ollabel{prop:total}
المجموعة $\fn{Tot}$ غير قابلة للحساب.
\end{prop}

\begin{proof}
نثبت انتفاء قابلية $\fn{Tot}$ للحساب باختزال $K$
إليها. فنعرّف $h(x,y)$ على الوجه الآتي:
\[
h(x,y) \simeq
\begin{cases}
0 & \text{إذا كان $x \in K$} \\
\fundefined & \text{خلاف ذلك.}
\end{cases}
\]
أما $h(x,y)$ فلا اعتماد لها على $y$ أصلًا.
ومن اليسير حساب $h$ جزئيًا: نتلقى $x,y$، ونحسب~$h$ بمحاكاة
الدالة~$\cfind{x}$ عند الدخل~$x$؛ فإن توقف الحساب، أعطت
$h(x,y)$ القيمة~$0$ ووقفت. فالدالة $h(x,y)$ هي
$\Zero(\umin{s}{T(x,x,s)})$، و$\Zero$ دالة الصفر الثابتة.

ومن مبرهنة $s$-$m$-$n$ نحصل على دالة عودية بدائية~$k(x)$، بحيث لكل
$x$ و~$y$،
\[
\cfind{k(x)}(y) \simeq
\begin{cases}
0 & \text{إذا كان $x \in K$} \\
\fundefined & \text{خلاف ذلك.}
\end{cases}
\]
فتكون $\cfind{k(x)}$ كلية حين $x \in K$، وغير معرّفة في أي موضع
إن لم يكن. فهذه $k$ اختزال من $K$ إلى~$\fn{Tot}$.
\end{proof}

\begin{digress}
بل إن $\fn{Tot}$ لا تقبل التعداد الحسابي أصلًا؛ فتعقيدها في
منزلة أعلى من «الهرم الحسابي». ولا نتناول هنا إثبات هذه الزيادة.
\end{digress}

\end{document}
